\documentclass[11pt]{article}

% Change "review" to "final" to generate the final (sometimes called camera-ready) version.
% Change to "preprint" to generate a non-anonymous version with page numbers.
%\usepackage[preprint]{acl}
\usepackage[final]{acl}

% Standard package includes
\usepackage{times}
\usepackage{latexsym}

% For proper rendering and hyphenation of words containing Latin characters (including in bib files)
\usepackage[T1]{fontenc}
% For Vietnamese characters
% \usepackage[T5]{fontenc}
% See https://www.latex-project.org/help/documentation/encguide.pdf for other character sets

% This assumes your files are encoded as UTF8
\usepackage[utf8]{inputenc}

% This is not strictly necessary, and may be commented out,
% but it will improve the layout of the manuscript,
% and will typically save some space.
\usepackage{microtype}

% This is also not strictly necessary, and may be commented out.
% However, it will improve the aesthetics of text in
% the typewriter font.
\usepackage{inconsolata}

%Including images in your LaTeX document requires adding
%additional package(s)
\usepackage{graphicx}

\usepackage{booktabs}
\usepackage{amsmath}

% If the title and author information does not fit in the area allocated, uncomment the following
%
%\setlength\titlebox{<dim>}
%
% and set <dim> to something 5cm or larger.

\title{Detecting Subtle Sense Shift with Polysemy-Aware Trends}

% Author information can be set in various styles:
% For several authors from the same institution:
% \author{Author 1 \and ... \and Author n \\
%         Address line \\ ... \\ Address line}
% if the names do not fit well on one line use
%         Author 1 \\ {\bf Author 2} \\ ... \\ {\bf Author n} \\
% For authors from different institutions:
% \author{Author 1 \\ Address line \\  ... \\ Address line
%         \And  ... \And
%         Author n \\ Address line \\ ... \\ Address line}
% To start a separate ``row'' of authors use \AND, as in
% \author{Author 1 \\ Address line \\  ... \\ Address line
%         \AND
%         Author 2 \\ Address line \\ ... \\ Address line \And
%         Author 3 \\ Address line \\ ... \\ Address line}
\iffalse
\author{Ondřej Herman \and Pavel Rychlý \\

  Natural Language Processing Centre \And Lexical Computing \\ 
  Masaryk University \\
  Brno, Czech Republic\\
  \texttt{\{xherman1,pary\}@fi.muni.cz}
}
\fi

\author{
\textbf{Ondřej Herman\textsuperscript{1,2}},
\textbf{Pavel Rychlý\textsuperscript{1,2}}
\\
\\
\texttt{\{xherman1,pary\}@fi.muni.cz}
\\
\\
\textsuperscript{1}Natural Language Processing Centre, Masaryk University, Brno, Czech Republic\\
\textsuperscript{2}Lexical Computing, Brno, Czech Republic
}

%\author{
%  \textbf{First Author\textsuperscript{1}},
%  \textbf{Second Author\textsuperscript{1,2}},
%  \textbf{Third T. Author\textsuperscript{1}},
%  \textbf{Fourth Author\textsuperscript{1}},
%\\
%  \textbf{Fifth Author\textsuperscript{1,2}},
%  \textbf{Sixth Author\textsuperscript{1}},
%  \textbf{Seventh Author\textsuperscript{1}},
%  \textbf{Eighth Author \textsuperscript{1,2,3,4}},
%\\
%  \textbf{Ninth Author\textsuperscript{1}},
%  \textbf{Tenth Author\textsuperscript{1}},
%  \textbf{Eleventh E. Author\textsuperscript{1,2,3,4,5}},
%  \textbf{Twelfth Author\textsuperscript{1}},
%\\
%  \textbf{Thirteenth Author\textsuperscript{3}},
%  \textbf{Fourteenth F. Author\textsuperscript{2,4}},
%  \textbf{Fifteenth Author\textsuperscript{1}},
%  \textbf{Sixteenth Author\textsuperscript{1}},
%\\
%  \textbf{Seventeenth S. Author\textsuperscript{4,5}},
%  \textbf{Eighteenth Author\textsuperscript{3,4}},
%  \textbf{Nineteenth N. Author\textsuperscript{2,5}},
%  \textbf{Twentieth Author\textsuperscript{1}}
%\\
%\\
%  \textsuperscript{1}Affiliation 1,
%  \textsuperscript{2}Affiliation 2,
%  \textsuperscript{3}Affiliation 3,
%  \textsuperscript{4}Affiliation 4,
%  \textsuperscript{5}Affiliation 5
%\\
%  \small{
%    \textbf{Correspondence:} \href{mailto:email@domain}{email@domain}
%  }
%}



\begin{document}
\maketitle
\begin{abstract}
\iffalse
We introduce a scalable, language‑independent pipeline that couples polysemy‑aware Adaptive Skip‑Gram induction with robust trend estimation to surface emerging word senses in continuously‑updated web corpora, enabling lexicographers to detect week‑to‑month semantic drifts—something traditional historical LSCD methods miss.
\fi

Language changes faster than dictionaries can be revised, yet automatic tools still struggle to spot the subtle, short‑term shifts in meaning that precede a formal update. We present a language‑independent pipeline that detects word‑sense shifts in large, time‑stamped web corpora. The method couples a robust re‑implementation of the Adaptive Skip‑Gram model, which induces multiple sense vectors per lemma without any external inventory, with a second stage that tracks each sense through time under three alternative frequency normalizations. Linear Regression and the robust Mann-Kendall/Theil–Sen estimator then test whether a sense’s frequency slope deviates significantly from zero, producing a ranked list of headwords whose semantics are drifting.
                                                                              
We evaluate the system on the English (12\,B tokens) and Czech (1\,B tokens) Timestamped corpora for May 2023–May 2025. Expert annotation of the top‑100 candidates for each model variant shows that 50.7\% of Czech and 25.7\% of English headwords exhibit genuine sense shifts, despite web‑scale noise.

\end{abstract}

\section{Introduction}
Dictionaries lag behind the language they purport to describe: while
web discourse spreads new shades of meaning within weeks,
lexicographic updates arrive months or years later.
Automating the early detection of \emph{word‑sense shifts}
would help editors allocate effort to the most dynamic headwords
and, in turn, keep lexical resources current.

Existing research on lexical semantic change
focuses on century‑scale, well‑curated corpora
and binary "changed / unchanged" judgements
\citep{schlechtweg2023human,tahmasebi2023computationalmodelingsemanticchange}.
We argue that today’s monitor corpora—
billions of time‑stamped web tokens—call for a
different objective: spotting \emph{emerging or drifting senses}
in near‑real time, despite noise and domain volatility.

We propose a three‑stage, language‑independent pipeline for tracking within-lemma sense redistribution that

\begin{enumerate}
  \item induces polysemy‑aware embeddings with an
        efficient \textbf{Adaptive Skip‑Gram} implementation (§\ref{sec:sense}),
  \item tracks each induced sense through time and tests
        for statistically significant trends (§\ref{sec:trend}),
  \item outputs ranked headwords that lexicographers can inspect (§\ref{sec:results}).
\end{enumerate}

Evaluated on \textbf{12\,B English} and \textbf{1\,B Czech} tokens (May 2023–May 2025),
the method delivers a \textbf{25–51\%} precision of genuinely new or rising senses,
according to expert annotation.





\section{Related Work}
\paragraph{Lexical Semantic Change Detection (LSCD).}
Early LSCD studies relied on co‑occurrence vectors
\citep{sagi2009semantic,gulordava2011distributional}
or topic models \citep{cook2013lexicographic}.
Recent competitions \citep{schlechtweg-etal-2020-semeval}
popularised embedding‑distance measures,
yet still treat time as two coarse epochs.

\paragraph{Sense‑aware embedding models.}
Multi‑prototype embeddings such as
Adaptive Skip‑Gram (Adagram) \citep{bartunov2016breaking}
and later transformer‑based approaches
capture polysemy more explicitly than single‑vector models.
However, few works couple them with diachronic trend analysis.

\paragraph{Trend detection in corpora.}
\citet{trendy} introduced frequency normalization and statistical trend estimation
for neologism detection.
We extend this idea to sense‑level frequencies.

\paragraph{Lexicographic applications.}
Commercial platforms (e.g. Sketch Engine; \citealp{kilgarriff2014sketch})
offer keyword‑in‑context views but no automatic sense‑shift alerts.
Our pipeline is designed to plug directly into such workflows.



\section{Data: Timestamped Web Corpora}
\label{sec:data}

We work with the English and Czech
Timestamped monitor corpora~\citep{trendcorpora2025} distributed in Sketch Engine.\footnote{Accessible through the interface at \url{https://www.sketchengine.eu} and made for linguistic research. The full corpus text is available for research purposes .} The corpus construction pipeline is described in~\citet{krek2017jsinewsfeed}.
Each document comes with an RSS publication timestamp.
The text is extracted from HTML using jusText~\citep{pomikalek2011removing}, tokenised by Unitok~\citep{michelfeit2014text}, and deduplicated on paragraph level using Onion~\citep{pomikalek2011removing}, and POS‑tagged + lemmatised using TreeTagger~\citep{schmid2013probabilistic}.

\begin{table}[t]
\centering
\small
\begin{tabular}{lrrr}
\toprule
Corpus & Tokens & Time span & Docs/day \\
\midrule
English & 12.0\,B & 2023‑05 – 2025‑05 & 21.1\,k \\
Czech & 1.1\,B & 2023‑05 – 2025‑05 & 2.6\,k \\
\bottomrule
\end{tabular}
\caption{Sketch Engine Monitor corpora used in this study.}
\end{table}

Their size and daily growth make them ideal test‑beds
for short‑term semantic drift.


\section{Methodology}
\label{sec:method}
Our pipeline decouples \emph{sense induction} (§\ref{sec:sense})
from \emph{trend estimation} (§\ref{sec:trend}),
each oblivious to language specific issues.
To identify trending senses, we build a word sense induction model based on the source corpus, and extract the senses of the examined headwords. We then calculate the frequencies of the senses over time, estimate the statistical trend, and rank the results.

\subsection{Word Sense Induction}
We use Adaptive Skip-gram~\citep{adagram} to induce word senses from the corpus text. This algorithm satisfies multiple practical constraints that arise when tracking short‑term meaning change in web‑scale monitor corpora:

\paragraph{Language and language-resource independence.} It learns multiple prototypes for a word directly from raw context, avoiding reliance on external sense inventories that either do not exist (for low‑resource languages) or lag behind emerging usage. This capability is essential for our goal of detecting senses that were \emph{unknown at training time} -- a scenario where large language models fall short, since their language knowledge is frozen at the moment of pre‑training.

\paragraph{Computationally efficient.} The training loop mirrors classic skip‑gram, with only a constant‑factor overhead for sense selection, enabling frequent retraining.

\paragraph{Representation quality.} Intrinsic evaluation on the ShadowSense benchmark~\citep{herman2024shadowsense} confirms that Adaptive skip-gram’s induced clusters align with human intuition for a majority of test words, providing the semantic fidelity required for downstream trend analysis.  While more recent transformer‑based multi‑prototype models achieve marginal gains on static datasets, their heavier resource demands and dependence on pre‑training make them ill‑suited for the continuously growing corpora that motivate our study. In short, Adagram delivers the right mix of quality, scalability, and resource independence needed to monitor lexical change in real time.

\subsection{Sense Induction with Adaptive Skip‑Gram}
\label{sec:sense}
To obtain the word sense induction model, we train Adaptive skip-gram on raw tokens, allowing up to $K{=}10$ senses per headword. To reach web scale, we re-implemented the original Julia code in Rust, leveraging compact Manatee~\citep{rychly2007manatee}
corpus representation for efficient access. The implementation is available at \url{https://github.com/ondra/sensetrends}.
Computing 64-dimensional embeddings over one training pass on 12\,B tokens
takes 60 hours using 32 threads on AMD EPYC 9654 CPU.\footnote{Using more threads does not speed up the computation significantly due to memory access contention.}

Each token is then assigned to its most probable sense and tabulated by the diachronic epoch.

\subsection{Diachronic Frequency Normalization}
\label{sec:freq}
At this point, we have obtained the frequency distributions of the senses over time for each word, which are not directly comparable. We explore three normalization approaches, each highlighting
a different aspect of change.

For every headword, $S$ is the set of senses and $E$ is the set of the epochs, (e.g. months) present in the corpus. Let $f_{\text{raw}}(s,e)$ be the count of sense $s \in S$
in epoch $e \in E$, $N(e)$ the token total of $e$.

\paragraph{Epoch‑normalized}  
      \[
        f_{en}(s,e)=\frac{|E|\;f_{\text{raw}}(s,e)/N(e)}
                         {\sum_{e'}f_{\text{raw}}(s,e')/N(e')}
      \]
      A natural extension of the approach used by~\citet{trendy} for the analysis of the diachronic behavior of words. Every sense is examined in isolation and sums to $|E|$, so for a perfectly stable sense, $f_{en}$ will be 1 everywhere. Peaks indicate epochs where $s$ is over-represented relative to its own history.

\paragraph{Global‑normalized}  
      \[
        f_{gn}(s,e)=\frac{|S||E|\;f_{\text{raw}}(s,e)/N(e)}
                         {\sum_{s'}\sum_{e'}f_{\text{raw}}(s',e')/N(e')}
      \]
      
      Similar to $f_{en}$, but reserves proportionality across the senses based on the corpus frequency. The dominance of rare senses is reduced compared to $f_{en}$, which treats every sense in isolation.

\paragraph{Sense‑relative}
      \[
        f_{sr}(s,e)=\frac{f_{\text{raw}}(s,e)}
                         {\sum_{s'}f_{\text{raw}}(s',e)}
      \]
      Captures \emph{within‑word} redistribution of meaning. The $f_{sr}$ sums to 1 over all senses within every epoch. This approach removes the dependence on absolute corpus frequencies and correlates with sense shift closely.

\subsection{Trend Estimation}
\label{sec:trend}
In this step, we estimate the rate of change and the corresponding $p$-value for every sense.

Given the time series $\{(x_i,y_i)\}$, where $x_i$ indexes epochs
and $y_i$ is the raw corpus frequency normalized by $f_{en}$, $f_{gn}$, or $f_{sr}$,
we fit (i) ordinary least squares (OLS) and
(ii) the Theil–Sen (TS) robust slope estimator accompanied by the Mann-Kendall statistical test as described in~\citet{trendy}.
In our experiment, a sense is flagged if \emph{p} < $10^{-4}$.

\subsection{Candidate Ranking}

For each headword, we keep the steepest trending sense,
then sort words by the slope magnitude.
The pipeline outputs top‑$k$ words
for review.

\section{Experimental Setup}
\paragraph{Vocabulary.}
We analyze the diachronic behavior of the $30,000$ most frequent lemmas,
excluding those starting with initial capital as a rudimentary form of named entity recognition.

\paragraph{Epoch size.}
Monthly windows offer a compromise between stability
and responsiveness ($\approx800$\,M English tokens per diachronic epoch).

\paragraph{Annotation protocol.}
A trained in-house linguist labeled the top 100
candidates sorted according to the trend slope from each of six method variants (3 norms and 2 slope estimators)
as \begin{enumerate}
    \item \textsc{OK} - potentially of lexicographical interest
    \item \textsc{Bad} - no lexicographical relevance / senses not understandable
    \item \textsc{Error} - other processing error / bad headword
\end{enumerate}

The total number of headwords was approximately 250 for each of the languages due to overlap between the lists.

The annotator was shown the frequency plot for every sense, with the senses described by a list of their nearest neighbors in the word embedding space of the Adagram model, as seen in Figure \ref{fig:rag}.

\section{Evaluation and Results}
\label{sec:results}


\begin{table}[h!]
\centering\small
\begin{tabular}{llccc}
\toprule
Norm & Est. & \multicolumn{3}{c}{English (12\,B)} \\
     &      & OK & Bad & Err \\
\midrule
$f_{en}$ & OLS & 28 & 57 & 15 \\
$f_{en}$ & TS  & 24 & 64 & 12 \\
$f_{gn}$ & OLS & 26 & 59 & 15 \\
$f_{gn}$ & TS  & 24 & 62 & 14 \\
$f_{sr}$ & OLS & 24 & 58 & 18 \\
$f_{sr}$ & TS  & 22 & 62 & 16 \\
\bottomrule
\end{tabular}
\caption{English precision at 100.}
\label{tab:en}
\end{table}

\begin{table}[h!]
\centering\small
\begin{tabular}{llccc}
\toprule
Norm & Est. & \multicolumn{3}{c}{Czech (1.1\,B)} \\
     &      & OK & Bad & Err \\
\midrule
$f_{en}$ & OLS & 49 & 47 & 4 \\
$f_{en}$ & TS  & 50 & 49 & 1 \\
$f_{gn}$ & OLS & 49 & 47 & 4 \\
$f_{gn}$ & TS  & 50 & 48 & 2 \\
$f_{sr}$ & OLS & 52 & 46 & 2 \\
$f_{sr}$ & TS  & 54 & 43 & 3 \\
\bottomrule
\end{tabular}
\caption{Czech precision at 100.}
\label{tab:cs}
\end{table}

\paragraph{Findings.}
The Czech corpus attains a mean 50.7\% hit‑rate,
twice the English 25.7\%.  We attribute the gap to heavier spam and boiler‑plate contamination in English web feeds, and larger absolute corpus size, which amplifies small, but still statistically significant artifacts.

\subsection{Sample results}
\subsubsection{Annotated as \textit{OK}}


\begin{figure*}[t]
\centering
\includegraphics[width=\linewidth]{rag-n.pdf}
\caption{Emerging AI sense of \emph{rag}.}
\label{fig:rag}
\end{figure*}

Figure \ref{fig:rag} plots weekly sense shares for the English noun \emph{rag};
a new AI‑related "retrieval-augmented generation" sense appears at the end of 2024.
Additional interesting examples include the English noun \emph{whale} (animal $\rightarrow$ influential investor), and Czech \emph{motorista} (motorist $\rightarrow$ member of a political party), or \emph{box} (trending sense describing a self-pickup delivery box).

\subsubsection{Annotated as \textit{Bad}}
The annotator understood the word, but the change was not meaningful, interesting, or the trending sense was not possible to distinguish from the others. For example, \textit{fertilizer} trending within stock news, or \textit{integration} trending in contexts related to market research, where the sense carried by the word is not truly different, but separated by the WSI algorithm due to its specificity, or \textit{signup} trending within fantasy footbal contexts, which were not present in earlier versions of the underlying corpus. 

\subsubsection{Annotated as \textit{Error}}
Typical representatives of tokens marked as \textit{error} are tokens, which were not recognized as words in the target language by the annotator, e.g. \textit{"6a"} trending in figure labels, \textit{+1} as a part of a phone number, \textit{\^} (a single caret character) trending within stock report contexts, or URLs or their parts.

\section{Discussion}
The expert evaluation found a \textbf{50.7\% precision} for Czech
versus \textbf{25.7\% for English}. While this might seem low, the evidence-based alternative is the direct analysis of the whole candidate list, so our method has the potential to save significant amount of work.

\subsection{Czech and English Result Difference}
\paragraph{Corpus composition.}
  The English Timestamped corpus is almost an order of magnitude larger.
  More sources imply more random shift bursts and boiler‑plate fragments that
  survive text cleaning and deduplication and inflate sense counts. English is an important language for our research, so we also tend to improve the processing and add new data sources more often compared to Czech, introducing shifts in frequencies.
\paragraph{Spam vulnerability.}  
  English is a primary target for SEO and click‑bait content that is hard
  to identify at crawl time; such noise produces statistically significant
  but linguistically meaningless trends.  Currently, websites appear to be all or nothing with respect to spam content, but the extracted text is locally plausible -- no boilerplate remains, and the issues only emerge in aggregate. LLMs in the hands of spammers significantly complicate the direct detection of spam text, so this issue will only get worse.
\paragraph{Relative domain stability.}
  Czech web feeds are fewer and more homogeneous, reducing abrupt
  register or genre shifts that can masquerade as semantic change.

\subsection{What about recall?}
Unfortunately, we are only able to report precision and not recall at this time, as we do not have any grounded data describing word sense shift along with a compatible corpus source.
As a proxy for estimating recall, we carried out experiments with sampling headwords, for which no significant trend was identified, but we failed to find any false negatives in a sample of 500 lemmas.
In the future, we intend to compare our result with a revision of a dictionary to obtain a more objective insight into the quality of the methods.

\subsection{Three Normalization Strategies}
The result lists for the three normalization strategies turned out to be very similar with a large overlap, contrary to our expectations. Even though the ordering of the result is not significantly different, the estimates have different and meaningful interpretations for the direct quantification of the change. Trend estimate of $f_{en}$ represents the diachronic change of a single sense in isolation; $f_{gn}$ introduces weighting by sense frequency within a single headword, and $f_{sr}$ quantifies the redistribution between different senses.

\section{Future Work}
\paragraph{Evaluation.} The current evaluation is quite limited, so we plan to explore the results for other languages and for larger result lists. We also intend to measure inter-annotator agreement.

\paragraph{Artifact cleanup.} The method is sensitive to changes in corpus composition shifts and to all kinds of processing changes and issues, such as encoding and boilerplate removal errors. These issues tend to be strongly associated with specific feeds, so we want to integrate feed-level weighting to these types of artifacts.

\paragraph{Improve sense descriptions.} Annotators found nearest neighbors cryptic or opaque. We believe tightly bound collocations as disambiguators should provide superior results.

\section{Conclusion}
We presented a pipeline that
combines polysemy‑aware embeddings with robust trend statistics
to identify short‑term word‑sense shifts in billion‑token web corpora.
The method achieves up to 54\% precision in Czech
and offers immediately actionable headword queues for lexicographers.
%Our open‑source release invites further cross‑lingual and methodological comparison.

\section*{Limitations}
\paragraph{Annotation coverage.}
Only one annotator judged the sample. Inter-annotator agreement remains to be measured.
\paragraph{Language coverage.}
We tested the results for two Indo-European languages. Other language families or low-resource languages may need different preprocessing.
\paragraph{Bias in web sources.}
Monitor corpora skew towards news outlets,
potentially missing sense change in spoken or social media registers.
\paragraph{New meaning or just different context?}
It is sometimes difficult to differentiate between a distinct sense of a word and a specific context a word appears in. WSI is a hard problem and remains unsolved~\citep{mosolova-etal-2025-llm}. Even agreement between humans is often low~\citep{kilgarriff1997don,herman2024shadowsense}.


\section*{Acknowledgments}
We used an AI assistant to improve the language of the article.

The work described herein has also been using tools provided by the LINDAT/CLARIAH-CZ Research Infrastructure (https://lindat.cz), supported by the Ministry of Education, Youth and Sports of the Czech Republic (Project No. LM2023062)

The described study comes from the project ``On our own: Opportunities and
Risks in the Individualization of Society (PRINS)
CZ.02.01.01/00/23\_025/0008710'', which is co-financed by the European
Union.

%We thank [ funding body ] for grant [ ID ],
%and \ldots. %the Sketch Engine team for access to the
%Timestamped corpora and infrastructure support.

% https://anonymous.4open.science/r/sensetrends-27BD/

% Bibliography entries for the entire Anthology, followed by custom entries
%\bibliography{anthology,custom}
% Custom bibliography entries only
\bibliography{custom}

\appendix

%\section{Example Appendix}
%\label{sec:appendix}
%This is an appendix.

\end{document}
